\begin{proposition}
  {\normalfont
  (Central Newton's formula \cite[p.~462]{butzer1989central}).}
  For a function $f$ defined on integers and arbitrary integers $x,t$
  \label{prop:newtons-interpolation-formula-in-central-differences}
  \begin{align*}
    f(x)
    = \sum_{k=0}^{\infty} \frac{\centralFactorial{(x-t)}{k}}{k!} \delta^{k} f(t),
  \end{align*}
  where $\delta^{k} f(t)$ denotes the $k$-th central finite difference
  \begin{align*}
    \delta^{k} f(t) = \tsum_{j=0}^{k} (-1)^j \tbinom{k}{j} f\!\left(t+\tfrac{k}{2} - j\right),
  \end{align*}
  and $\centralFactorial{(x-t)}{k}$ are central factorials
  defined by~\eqref{def:central-factorials}.
\end{proposition}


\begin{lemma}[Newton's formula for powers]
  \label{prop:newtons-formula-for-powers}
  For integers $n,m \geq 0$, and an arbitrary integer $t$,
  \begin{align*}
    n^m = \sum_{k=0}^{m} \frac{\centralFactorial{(n-t)}{k}}{k!} \delta^{k} t^m.
  \end{align*}
\end{lemma}

\begin{lemma}[Central factorials binomial form]
  \label{lem:binomial-form-of-central-factorials}
  For integers $n$ and $k \geq 1$,
  \begin{align*}
    \frac{\centralFactorial{n}{k}}{k!}
    = \frac{n}{k} \binom{n+\frac{k}{2}-1}{k-1}.
  \end{align*}
  \begin{proof}
    We have
    \begin{align*}
      \frac{\centralFactorial{n}{k}}{k!}
      &= \frac{n}{k!}\fallingFactorial{n+\frac{k}{2}-1}{k-1}
      = \frac{n}{k (k-1)!} \fallingFactorial{n+\frac{k}{2}-1}{k-1}
      = \frac{n}{k} \binom{n+\frac{k}{2}-1}{k-1},
    \end{align*}
    because of the identity in falling factorials
    $\frac{\fallingFactorial{x}{n}}{n!} = \binom{x}{n}$.
  \end{proof}
\end{lemma}

\begin{lemma}[Central factorials binomial form decomposition]
  \label{lem:halved-central-factorials}
  For integers $n,k\geq 0$,
  \begin{align*}
    \frac{\centralFactorial{n}{k}}{k!}
    = \frac{n}{k} \binom{n+\frac{k}{2}-1}{k-1}
    = \frac{1}{2} \binom{n+\frac{k}{2}}{k} + \frac{1}{2} \binom{n+\frac{k}{2}-1}{k}.
  \end{align*}
  \begin{proof}
    We have $\binom{a}{k} = \frac{a}{k} \binom{a-1}{k-1}$.
    Let $a=n + \frac{k}{2}$, then
    \begin{align*}
      \tbinom{n + \frac{k}{2}}{k} = \tfrac{n + \tfrac{k}{2}}{k} \tbinom{n+\frac{k}{2}-1}{k-1}
      = \tfrac{n}{k} \tbinom{n+\frac{k}{2}-1}{k-1} + \tfrac{1}{2} \tbinom{n+\frac{k}{2}-1}{k-1}.
    \end{align*}
    Thus,
    \begin{align*}
      \tfrac{n}{k} \tbinom{n+\frac{k}{2}-1}{k-1} = \tbinom{n + \frac{k}{2}}{k} - \tfrac{1}{2} \tbinom{n+\frac{k}{2}-1}{k-1}.
    \end{align*}
    By moving under common denominator, and applying binomial recurrence yields
    \begin{align*}
      \tfrac{n}{k} \tbinom{n+\frac{k}{2}-1}{k-1}
      = \tfrac{1}{2}
      \left[
        2 \tbinom{n + \frac{k}{2}}{k} - \tbinom{n+\frac{k}{2}-1}{k-1}
      \right]
      = \tfrac{1}{2}
      \left[
        \tbinom{n + \frac{k}{2}}{k} + \tbinom{n + \frac{k}{2}-1}{k} + \tbinom{n + \frac{k}{2}-1}{k-1} - \tbinom{n+\frac{k}{2}-1}{k-1}
      \right].
    \end{align*}
    Thus, the statement follows.
  \end{proof}
\end{lemma}

\begin{lemma}[Newton's formula for powers decomposition]
  \label{prop:halved-central-factorial-powers}
  For integers $n,m \geq 0$, and an arbitrary integer $t$,
  \begin{align*}
    n^m = \frac{1}{2} \sum_{k=0}^{m}
    \left[
      \binom{n-t+\frac{k}{2}}{k} + \binom{n-t+\frac{k}{2}-1}{k}
    \right] \delta^{k} t^m.
  \end{align*}
\end{lemma}

\begin{lemma}[Generalized hockey-stick identity]
  \label{prop:generalized-hockey-stick-identity}
  For integers $a,b$ and $j$,
  \begin{align*}
    \sum_{k=a}^{b} \binom{k}{j} = \binom{b+1}{j+1} - \binom{a}{j+1}.
  \end{align*}
  \begin{proof}
    We have,
    $\sum_{k=a}^{b} \binom{k}{j}
    = \binom{a}{j} + \binom{a+1}{j} + \cdots + \binom{b}{j}
    = \left( \sum_{k=0}^{b} \binom{k}{j} \right) - \left( \sum_{k=0}^{a-1} \binom{k}{j} \right)$.
    By hockey stick identity $\sum_{k=0}^{n} \binom{k}{j} = \binom{n+1}{j+1}$ yields,
    \begin{align*}
      \tsum_{k=a}^{b} \tbinom{k}{j}
      = \left( \tsum_{k=0}^{b} \tbinom{k}{j} \right) - \left( \tsum_{k=0}^{a-1} \tbinom{k}{j} \right)
      = \tbinom{b+1}{j+1} - \tbinom{a}{j+1}.
    \end{align*}
    This completes the proof.
  \end{proof}
\end{lemma}

\begin{lemma}[Centered hockey-stick identity]
  \label{lem:centered-hockey-stick-identity}
  For integers $n\geq1$, and arbitrary integers $t,k,r$
  \begin{align*}
    \sum_{j=1}^{n} \binom{j -t + \frac{k}{2} +r}{k}
    = \binom{n -t + \frac{k}{2} + r +1}{k+1} - \binom{1 -t + \tfrac{k}{2} +r}{k+1}.
  \end{align*}
  \begin{proof}
    By setting $j=1 -t + \tfrac{k}{2} +r$, and $b=n-t+ \tfrac{k}{2} +r$ yields
    \begin{align*}
      \tsum_{j=1}^{n} \tbinom{j -t + \tfrac{k}{2} +r}{k} =
      \tsum_{j=1 -t + \tfrac{k}{2} +r}^{n -t + \tfrac{k}{2} +r} \tbinom{j}{k}.
    \end{align*}
    By generalized hockey-stick identity~\eqref{prop:generalized-hockey-stick-identity},
    \begin{align*}
      \tsum_{j=1 -t + \tfrac{k}{2} +r}^{n -t + \tfrac{k}{2} +r} \tbinom{j}{k}
      = \tbinom{n -t + \tfrac{k}{2} + r +1}{k+1} - \tbinom{1 -t + \tfrac{k}{2} +r}{k+1}.
    \end{align*}
    Thus, the claim follows.
  \end{proof}
\end{lemma}
